\sectionmark{Chemolithotrophie}

\section{Chemolithotrophie}
\begin{enumerate}
	\item Wie konnte Winogradsky zeigen, dass die Schwefelbakterien einen chemolithotrophen Stoffwechsel haben?
		
		Durch Beobachtung an einem Bachlauf mit hohem \ce{H2S}-Gehalt.
		Er stellte fest,
		dass bei höheren Konzentrationen von Schwefelwasserstoff auch mehr Schwefel einlagerten.
		So Schlussfolgertee er,
		dass die Bakterien \ce{H2S} oxidierten und als molkeularen Schwefel ablagerten
		und das Mikroorganismen auch in der Lage sind anorganische Substanzen zu oxidieren.

		Chemolithotrophie wird auch als Chemoautotrophie bezeichnet.

	\item Erklären Sie anhand der Redoxpotentiale der Reaktionen, 
			welche besonderen Anforderungen die chemolithotrophe Lebensweise erfordert, 
			die bei heterotrophen Wachstum auf Kohlenhydraten nicht auftreten?

		\begin{equation}
			2\ \ce{H} + \ce{O2} \rightarrow 2\ \ce{H2O} \hspace{1.8cm} \Delta G^0 = -278 kJ
			\label{Hydrolyse}
		\end{equation}

		Die in Formel \ref{Hydrolyse} beschriebene Redox-Reaktion erzeugt.
		\ce{NADH} wird bei Chemolithotrophen Bakterien die diese Reation betreiben
		auch mit einer nicht Membran gebundenen Hydrogenase regeregeneriert.


	\item Welche Wege der Oxidation reduzierter Schwefelverbindungen kennen Sie?
		
		\begin{equation}
			\ce{H2S} + 2\ \ce{H2O} \rightarrow \ce{SO4}^{2-} + 2\ \ce{H+}
			\label{schwefelred}
		\end{equation}

		In Formel \ref{schwefelred} ist der Grundlegende Reaktion dargestellt.
		Diese Reaktion kann über verschiedene Wege realisiert sein.
		Für die Substratketten-Phosphorylierung und den Sulfit-Oxidase-Weg wird der Schwefel-Wasserstoff
		zunächst in \ce{SO3}$^{-2}$ überführt.
		Ausgangspunkte können deshalb auch elementarer Schwefel oder Thiosulfat (\ce{S2O3}$^{-2})$) sein.
		Die dabei endstehenden Elteronen können über das Elektronen-Transport-System weiter verwendet werden.

		\begin{description}
			\item[SOX-Weg] \hfill \\
				Neben der Sulfit-Oxidase-abhängigen Reaktion,
				kann Sulfid von vielen phototrophen Schwefelbakterien (z.B. \emph{Paracoccus pantotrophus})
				direkt über das sogenannte Sox-System oxidiert werden.
				Dieses System stellt die nötige Reduktionskraft zur Kohlenstoffdioxid-Fixierung bereit.
			\item[Substratketten-Phosphorylierung (APS-Weg)] \hfill \\
				Die Umsetzung von Sulfit zu Sulfat kann bei einigen \emph{Thiobacilli} mit einer
				Substratkettenphosphrylierung erfolgen.
				Dabei wird von einer Adenosinphosphosulfat-Reductase (APS-Rductase) AMP zu ADP umgesetzt.
				Eine Kinase erzeugt aus dem ADP wiederum ATP.

				Die APS-Reductase arbeitet gegenläufig der dissimilatorischen Sulfatredukion,
				siehe Abschnitt \ref{sec:anaerobeAtmung} "Anaerobe Atmung".
				
			\item[Sulfit-Oxidase-Weg] \hfill \\
				Unter Abspaltung von zwei Elektronen von Sulfit,
				endsteht Sulfat.
				\begin{equation}
					Sulfit-Oxidase(\ce{SO3}^{2-}) \rightarrow \ce{SO4}^{2-} + 2 e^-
					\label{sulfitoxidase}
				\end{equation}
		\end{description}

	\item Welche Aufgabe haben Nitrifizierer und Nitrosifizierer in der Abwasserbehandlung?
		
		Die Umsetzung von Ammonium (\ce{NH3}) in Nitrat und Nitrit.
		Dieser Vorgang lässt sich zweiteilen (Nitrifizierer/Nitrosifizierer). %TODO Nachlesen

		Der Abbau bestimmter Schwefelhaltiger Substanzen ist zu dem erwünscht,
		da diese zum Teil toxische sind (\ce{H2S}).

	\item In einem geschichteten Sediment laufen die Prozesse der anaeroben Atmung und der
			Oxidation reduzierter Verbindungen durch Chemolithotrophe ab. Ordnen Sie die
			Prozesse den Schichten zu und begründen Sie die Zuordnung!

		Chemolithotrophe Mikroorganismen siedeln sich,
		von oben betrachtet,
		frühstens in der Schichthöhe an in der nicht mehr aussreichend Licht für eine
		Phtotrophe Lebensweise vorhanden ist.

		Hier befindet sich auch die Grenzschicht,
		in der noch genug Sauerstoff vorhanden ist,
		so dass dieser als Elektronen-Akzeptor verwendet werden kann.
		Hier finden sich Nitrifizierern und Nitrosifizierer. %TODO RLY?

		Steht kein Sauerstoff mehr zur Verfügung,
		wird beispielsweise Ammonium und Nitrit zu Stickstoff und Sauerstoff oxidiert
		wie in Formel \ref{AnaerobeAmmoniumOxidation} dargestellt.

		\begin{equation}
			\ce{NH4}^+ + \ce{NO2}^- \rightarrow \ce{N2} + 2 \ce{H2O}
			\label{AnaerobeAmmoniumOxidation}
		\end{equation}

		Dies kann jedoch nur in einem abgeschlossenen Kompartiementen
		dem Anammoxosome,
		geschehen,
		da als Zwischenprodukt toxisches Hydrazin (\ce{N2H4}) endsteht.
		Man vermutet,
		dass die ATP Erzeugung hierbei über einen Protonengradient geschieht.

	\item Nehmen Sie einen Stammbaum der Bacteria und markieren Sie die chemolithotrophen Organismengruppen. 
		
		\begin{table}[h!]
		\begin{center}
		\begin{tabular}{l l} 
		\toprule
		Gattung			&	Eigenschaften 					\\
		\midrule
		Aquifex			&	fakultativer Anaerobier		\\
		Planctomyceten	&	anaerobe Ammoniumoxidaton	\\
		Deferribacter	&	\ce{H} \textrightarrow \ce{S}/\ce{NO3}$^-$ \\
		Nitrospira		&	Nitrifizierer, ausserhalb der Proteobacter \\
		Proteobacter	&	Schwefel- und Eisen-Oxidierer, Nitro- und Nitrifizierer \\
		\bottomrule
		\end{tabular}
		\caption{Chemolithotrophe Bacteria.}
		\label{tab:chemolithoBac}
		\end{center}
		\end{table}

		Chemolithotrophie ist eine ursprüngliche Eigenschaft der Mikroorganismen,
		weshalb wahrscheinlich auch die Ersten chemolitotoph waren.

	\item Welche Gruppen der Bacteria betreiben oxygene bzw. anoxygene Photosynthese?
		
		\begin{table}[h!]
		\begin{center}
		\begin{tabular}{l l} 
		\toprule	
		Gattung			&	Eigenschaften 						\\
		\midrule
		\multicolumn{2}{l}{Oxygene Photosynthese}			\\
		\midrule
		\multicolumn{2}{l}{Anoxygene Photosynthese}		\\
		Grüne Nichtschwefelbakterien 	& 	\\
		Grüne Schwefelbakterien			&	\\
		Gruppen der Proteobacteria		&	\\
		\bottomrule
		\end{tabular}
		\caption{Oxygene und Anoxygene Photosynthese.}
		\label{tab:AnOxPs}
		\end{center}
		\end{table}

		Oxygene Photosynthese wird nur von Cyanobakterien und den grünen Pflanzen betrieben.
		Alle anderen phototrophen Organismen betreiben eine anoxygene Photosynthese.

	
	\item Nenn Sie die wesentlichen Unterschiede zwischen den Prozessen der oxygenen und anoxygenen Photosynthese!

		Bei der Oxygenen Photosynthese wird das Licht sowohl für die Herstellung von \ce{ATP},
		als auch für die Herstellung von Reduktionsäquivalenten (\ce{NADH} und \ce{NADPH}).
		In der oxygenen Photosynthese dient ausschließlich Wasser als Elektronendonnator,
		für die Erzeugung von Reduktionsäquivalenten.
		Ein anderer Hauptunterschied ist die Namensgebende Erzeugung von elementarem
		Sauerstoff als Abfallprodukt der Oxygenen Photosynthese.
	
		Physiologisch unterscheiden sich die Photosynthesetypen durch die Verwendung
		verschiedener Pigmente.
		Bei der oxygenen Photosynthese kommen Chlorophylle zum Einsatz,
		bei der anoxygenen Photosynthese sind es Bakteriochorophylle.

	\item Wie wird bei der Photosynthese ein Protonengradient erzeugt?

		Durch den Elektronenfluss durch die Membran.
		Dabei werden vom Cytochromkomplex 4 \ce{H+} ans Stroma und 8 \ce{H+} ins Lumen abegegeben.

	\item Welche akzessorischen Pigmente kennen Sie? 

		\begin{description}
			\item[Carotinoide]
				sind die am weitestenverbreiteten akzessorischen Pigmente.
				Ihre Hautpaufgabe besteht im Schutz vor und der Aufnahme von Licht.
			\item[Phycobiline] 
				sind die für die ``Lichternte'' zuständige Proteine (Hauptantennenpigmente).
				Die Proteinen treten in großen Phycobilisomen auf und leiten letztlich die
				eingefangene Lichtenergie an das Reaktionszentrum weiter.
		\end{description}

		Durch die akzessorischen Proteine wird das Spektrum des absorbierten Lichtes vergrößert
		und somit kann mehr Energie ins Reaktive Zentrom des Chorophylls geleitet werden.

	\item In einem Teich sind die unterschiedlichen photosynthetischen Organismen nicht gleichmäßig verteilt. Warum?

		Es gibt einge große Diversität bei den Absorptionspektren der Photosysteme
		und der akzessorischen Pigemente.
		Dadurch verteilen sich die Organismen in die ökologischen Nischen
		in denen der Konkurrenzdruck nicht zu groß ist
		und Licht in den entsprechenden Wellenlängen vorhanden ist.

		So ergibt sich eine höhen Staffelung innerhalb von Gewässern.

	\item Wie wird Kohlendioxid fixiert?

		Die Autotrophe \ce{CO2}-Fixierung zumeist mit dem Calvin-Zyklus.
		
		\begin{figure}[ht!]
		\leavevmode
		\begin{center}
		\includegraphics[scale=0.33]{./pictures/calvin_simple_1000}
		\end{center}
		\caption{Einfache Darstellung des Calvin Zyklus mit seinen drei Phase. 
		Abbildung \ref{fig:calvinDetail} in Abschnitt \ref{sec:facts} Begriffe,
		zeigt eine genauere Darstellung.}
		\label{fig:calvinSimple}
		% http://de.wikipedia.org/w/index.php?title=Datei:Calvin-Zyklus_de.svg&filetimestamp=20090601125806
		\end{figure}

		Der Kreislauf erfordert \ce{NAD(P)H} und \ce{ATP} sowie die Schlüsselenzyme
		\slshape{Ribulosebisphosphat-Carboxylase} (RUBISCO)
		und \slshape{Phosphoribulose-Kinase} \cite[S.657]{brock}.
		In Formel \ref{calvin} ist die Gesamstöchioetrie des Zyklus dargestellt.
		
		\begin{equation}
			6 \ce{CO2} + 12 \ce{NADPH} + 18 \ce{ATP} \rightarrow \ce{C6H12O6(PO3H2)} + 12 \ce{NADP+} + 17 \ce{P}_a
			\label{calvin}
		\end{equation}

		Weiter Varianten der \ce{CO2}-Fixierung sind der umgekehrte Citrat-Zyklus bei den Grünen-Schwelfelbakterien findet,
		sowie der Hydroxypropionatweg der Grünen-Nichtschwefelbakterien.

\end{enumerate}
